\documentclass[10pt,a4paper]{amsart}
\usepackage[margin=19mm]{geometry}
\usepackage{amsmath,amssymb,mathtools}
\usepackage[scr=rsfs,cal=euler]{mathalfa}
\usepackage{xfrac}
\usepackage[unicode,hidelinks,bookmarks=false]{hyperref}
\newcommand{\ie}{i.e.\ }
\newcommand{\cf}{cf.\ }
\newcommand{\resp}{respectively\ }
\newcommand{\Subsection}{\subsection}
\providecommand{\nobreakdash}{\nobreak\hbox{-}}
\setlength{\emergencystretch}{2em}
\multlinegap0pt
\usepackage{siunitx}
\usepackage{siunitx}
\usepackage{siunitx}
\usepackage{siunitx}
\usepackage{bm}
\usepackage{amssymb}
\usepackage{enumerate}
\usepackage{siunitx}
\usepackage[shortlabels]{enumitem}
\usepackage{tikz}
\newtheorem{prop}{Proposition}
\newtheorem{lem}{Lemma}
\newcommand{\bsigma}{\bm{\sigma}}
\newcommand{\e}{\bm{e}}
\title{The search for alternating surgeries}
\author{Kenneth L. Baker, Marc Kegel, Duncan McCoy}
\date{Source: arXiv:2409.09842v2 (CC BY 4.0)}
\begin{document}
\maketitle

\noindent\emph{Source and licence.} The search for alternating surgeries. Version arXiv:2409.09842v2 (CC BY 4.0), distributed by arXiv under the Creative Commons Attribution 4.0 International licence, http://creativecommons.org/licenses/by/4.0/, as recorded in the arXiv licence metadata for this exact version and shown on the abstract page https://arxiv.org/abs/2409.09842v2. Full licence notice: \texttt{licenses/arxiv-2409.09842-CC-BY-4.0.txt}.\par\smallskip

\noindent\textit{Editorial scope.} Counted unit: the article's lem lem:sigma\_m\_bound (source0.tex lines 1357--1359). The article's complete original proof of it is delivered with it (source0.tex lines 1360--1441, one \texttt{proof} environment of 82 lines). No mathematics is written, repaired or re-derived: the statement and the proof are the article's own lines, in the article's own notation, and no retained line is substituted or re-flowed.  Retained uncounted as the source-local context the proof needs: lines 559--567.  Nothing was omitted from any retained span, so the package carries no omission record.  No retained line contains a citation, so the wrapper carries no bibliography and no bibliography entry.  No retained line contains a figure environment, an image-inclusion command or drawing code, so no image is delivered and the package declares no asset dependency.  Editorial formatting only, in addition: an AMS \texttt{amsart} preamble that restates the article's own declarations and macros used by the retained lines, namely the environments \texttt{prop}, \texttt{lem} on separate counters, together with the standard packages those lines need.  The retained environments print in this document as Proposition 1, Lemma 1 in order of appearance, whereas the article prints the same items with the numbering of its own full document.  The complete native source of the article as distributed in the arXiv e-print for this exact version is bundled unchanged as \texttt{sources/165299/source0.tex}.
\begin{prop}\label{prop:CM_condition}
Let $\bsigma = \sigma_1 \e_1 + \dots + \sigma_r \e_r$ be a changemaker vector.
Then for any $k\leq r$ and any integer $T\leq \sigma_1 + \dots + \sigma_k$, there is a subset $A\subseteq \{1,\dots, k\}$ such that 
\[
\pushQED{\qed} 
T=\sum_{i\in A}\sigma_i.\qedhere
\popQED
\]
\end{prop}

\begin{lem}\label{lem:sigma_m_bound}
If $\bm{z}\in L$ is irreducible, then $|z_m|\leq \sigma_m$.
\end{lem}

\begin{proof}
Without loss of generality let us suppose that $z_m>0$. Consider also the sets
\[
A_+=\{i\,|\, z_i\geq 1, 1\leq i \leq m-1\}
\]
and
\[
A_-=\{i\,|\, z_i\leq -1, 1\leq i \leq m-1\}.
\]
First we observe that at least one of $A_+$ or $A_-$ must be empty. If we had $a\in A_+$ and $b\in A_-$, then $\bm{x}=\e_a-\e_b\in L$ would be a vector $\bm{x}\neq \bm{z}$ satisfying 
\[
\bm{x}\cdot (\bm{z}-\bm{x})=z_a-z_b-2\geq 0.
\]
Since this contradicts the irreducibility of $\bm{z}$, we see that at least one of $A_+$ or $A_-$ must be empty.

Note that if $\sigma_m$ is tight, then $\sigma_m=m$ and if $\sigma_m$ is not tight, then $\sigma_m\leq m-1$. We prove the lemma by considering three separate cases:
\begin{enumerate}%[(a)]
    \item\label{it:A-non-empty} $A_-$ is non-empty and $\sigma_m$ is tight.
    \item\label{it:A+small} $\sigma_m$ is not tight and $|A_+|+\sigma_m\leq m-1$
    \item\label{it:A+large} $A_-$ is empty and $|A_+|+\sigma_m\geq m$.
\end{enumerate}
We note that these three cases encompass all possibilities. Together \eqref{it:A-non-empty} and \eqref{it:A+small} cover the case that $A_-$ is non-empty. The case that $A_-$ is empty is covered by \eqref{it:A+small} and \eqref{it:A+large}. In particular, note that if $\sigma_m$ is tight, then the condition $|A_+|+\sigma_m\geq m$ automatically holds. 

First, suppose that \eqref{it:A-non-empty} holds: we have $A_-$ is non-empty and $\sigma_m=m$. Without loss of generality we can assume that $1\in A_-$, i.e we have $z_1\leq -1$ and $z_i\leq 0$ for $i=2,\dots, m-1$. Consider $\bm{x}=\e_m-\e_{m-1}-\dots - \e_2 -2\e_1\in L$. Note that if $\bm{z}=\bm{x}$, then $z_m=1$. Thus we can assume that $\bm{x}\neq \bm{z}$.
The irreducibility of $\bm{z}$ gives 
\begin{align*}
    -1&\geq \bm{x}\cdot (\bm{z}-\bm{x}) \\
    &= z_m-1 - \sum_{i=2}^{m-1} (1+z_i) -2(z_1+2)\\
    & \geq z_m -1 -(m-2) -2\\
    &= z_m -1 -m.
\end{align*}
Rearranging gives $z_m\leq m=\sigma_m$, which is the desired bound if $A_-$ is nonempty and $\sigma_m$ is tight.

Now suppose that \eqref{it:A+small} holds: we have $\sigma_m\leq m-1$, and that $|A_+|+\sigma_m\leq m-1$. Note that this encompasses the case that $A_+$ is empty. In this case we consider $\bm{x}=\e_m-\e_1-\dots - \e_{\sigma_m}\in L$. The hypothesis that $|A_+|+\sigma_m\leq m-1$ implies that after relabelling the $\e_i$ we can assume that $A_+ \cap \{1, \dots, \sigma_m\}=\emptyset$, that is that $z_i\leq 0$ for $i=1, \dots, \sigma_m$. Again we can assume that $\bm{z}\neq \bm{x}$ since otherwise we would have $z_m=1$. Thus we find that
\begin{align*}
    -1&\geq \bm{x}\cdot (\bm{z}-\bm{x}) \\
    &= z_m-1 - \sum_{i=1}^{\sigma_m} (z_i+1)\\
    &\geq z_m -1 -\sigma_m.
\end{align*}
Rearranging this gives $z_m\leq \sigma_m$. Thus we have established the required bound when $|A_+|+\sigma_m\leq m-1$.

Finally suppose that \eqref{it:A+large} holds: we have that $A_-$ is empty and $|A_+|+\sigma_m\geq m$. Under these conditions we have that $z_i\geq 0$ for $i=1,\dots, m-1$ and, by assumption, that $z_m>0$. Since every non-zero vector in a changemaker lattice must have non-zero coordinates of both signs, there is some $i$ with $z_i<0$. Moreover we see that such an $i$ satisfies $i>m$. We take $g>m$ to be minimal such that $z_g\leq 0$.

By Proposition~\ref{prop:CM_condition}, there is a subset
$B\subset \{1, \dots, g-1 \}$ such that
\[
\sigma_g - \sigma_m= \sum_{i\in B}\sigma_i.
\]
Thus we can define an element $\bm{x}\in L$ of the form
\[
\bm{x}=-\e_g + \varepsilon \e_m + \sum_{i\in C} \e_i +\sum_{j\in D} \e_j,
\]
where $C=\{1,\dots, m-1\}\cap B$, $D= \{m+1, \dots, g-1\}\cap B$ and
\[
\varepsilon=\begin{cases}
    1 &\text{if $m\not\in B$}\\
    2 &\text{if $m\in B$.}
\end{cases}
\]
Furthermore, we assume that we have chosen $B$ so that $A_+\cap C$ is as large as possible. I.e so that $|C\setminus A_+|=\max\{0, |C|-|A_+| \}$. If $\bm{z}=\bm{x}$, then $z_m=\varepsilon\leq 2$, which is the desired bound. Thus we assume that $\bm{z}\neq \bm{x}$.
Since $\bm{z}$ is assumed to be irreducible, we have 
\begin{align*}
    -1&\geq \bm{x}\cdot (\bm{z}-\bm{x}) \\
    &=-z_g-1 + \varepsilon(z_m -\varepsilon) + \sum_{i\in C} (z_i-1) +\sum_{j\in D} (z_j-1)\\
    &\geq -1 + \varepsilon(z_m -\varepsilon) + \sum_{i\in C} (z_i-1)\\
    &\geq -1 + \varepsilon(z_m -\varepsilon) -|C\setminus A_+|.
\end{align*}
Thus rearranging we get the bound
\[
z_m\leq \frac{|C\setminus A_+|}{\varepsilon} +\varepsilon.
\]
Firstly, note that if $|C\setminus A_+|=0$, then we get the bound $z_m\leq \varepsilon\leq 2$. Thus we can assume that
$|C\setminus A_+|>0$. Since $|C|\leq m-1$ and $|A_+|+\sigma_m\geq m$, we have 
\[
|C\setminus A_+|=|C|-|A_+|\leq m-1-(m-\sigma_m)\leq \sigma_m -1.
\]
Since $\varepsilon\in \{1,2\}$, this gives the bound
\[
z_m\leq \max\left\{\sigma_m, \frac{\sigma_m+3}{2}\right\}.
\]
Since $\sigma_m\geq 2$ and $z_m$ is an integer, this is sufficient to show $z_m \leq \sigma_m$. This concludes the final case of the lemma.
\end{proof}

\end{document}
